\documentclass[a4paper,12pt]{article}

\usepackage[utf8]{inputenc}
\usepackage[T1]{fontenc}
\usepackage[table]{xcolor}
\usepackage{geometry}
\usepackage{longtable}
\usepackage{enumitem}
\usepackage{hyperref}
\usepackage{listings}
\usepackage{titlesec}
\usepackage{booktabs}
\usepackage{array}

\geometry{margin=2.5cm}

\definecolor{criticalcolor}{HTML}{DC3545}
\definecolor{highcolor}{HTML}{FD7E14}
\definecolor{mediumcolor}{HTML}{FFC107}
\definecolor{lowcolor}{HTML}{28A745}
\definecolor{resolvedcolor}{HTML}{6C757D}

\newcommand{\severitybox}[2]{%
  \colorbox{#1}{\makebox[2.5cm][c]{\textcolor{white}{\textbf{#2}}}}%
}

\lstset{
  basicstyle=\small\ttfamily,
  breaklines=true,
  frame=single,
  numbers=none,
  tabsize=2,
  backgroundcolor=\color{gray!10},
  extendedchars=true,
  inputencoding=utf8,
  literate={—}{---}1 {·}{$\cdot$}1 {≈}{$\approx$}1 {×}{$\times$}1
}

\title{\textbf{Performance Report \#10}\\[4pt]
\large Vulkan Engine — Shader Performance Audit (Round 2)}
\author{Henrique Rocha}
\date{\today}

\begin{document}
\maketitle

\begin{abstract}
This is the second dedicated audit of the GPU shader tree
(\texttt{shaders/*.\{vert,frag,geom,tesc,tese,comp\}} and
\texttt{shaders/includes/*.glsl}). It was produced from a clean re-read of
every shader, deliberately ignoring the 12 issues already covered in Report \#09
(which the user has since fixed, except for the two debug-only items). The goal
is identical: \emph{performance only}. No finding below changes the rendered
result — all features (tessellated water, triplanar terrain, EVSM shadows, GPU
vegetation, impostors, procedural sky) are preserved. 6 new issues are reported
(1 high, 4 medium, 1 low). The single high-severity item is a compute-shader
workgroup-size misconfiguration that leaves the GPU almost completely idle during
vegetation instance generation; the remaining items are redundant
per-invocation work that can be hoisted or merged.
\end{abstract}

\tableofcontents
\newpage

\section{Scope \& Methodology}

The audit covers the full shader tree plus the shared GLSL includes. Every file
was read in full and analysed for per-invocation cost, with emphasis on stages
that run at high invocation counts (compute dispatch, tessellation evaluation,
and the impostor geometry shaders). For each candidate I verified that the
redundant work is genuinely invariant (constant across invocations of the same
dispatch/draw) or recomputed from identical inputs, so removing it cannot alter
the output. As with Report \#09, the constraint from \texttt{AGENTS.md} is
respected: functionality is kept; only provably redundant or hoistable work is
targeted.

\subsection*{Shader Hotspot Overview}

\begin{center}
\begin{tabular}{lcl}
\toprule
\textbf{Shader / Include} & \textbf{Lines} & \textbf{Dominant cost} \\
\midrule
\texttt{vegetation\_instance\_gen.comp} & 82  & \texttt{local\_size\_x = 1} (serial) \\
\texttt{indirect.comp}                  & 117 & 6 plane normalizations / invocation \\
\texttt{vegetation\_cull.comp}          & 91  & 6 plane normalizations / invocation \\
\texttt{includes/perlin2d.glsl}        & 60  & \texttt{cos/sin} gradient hash \\
\texttt{includes/displacement.glsl}    & 50  & 3× triplanar-weight recompute \\
\texttt{includes/common.glsl}          & 58  & \texttt{sampleHeightTriplanar} redundancy \\
\texttt{impostors.geom}                & 153 & 20× Fibonacci dir build / primitive \\
\texttt{impostors\_depth.geom}         & 116 & 20× Fibonacci dir build / primitive \\
\texttt{vegetation.geom}               & 187 & 20× Fibonacci dir build / primitive \\
\texttt{sky.frag} / \texttt{sky\_equirect.frag} & 96 / 120 & residual \texttt{sin} hash \\
\bottomrule
\end{tabular}
\end{center}

\newpage

\section{Issue Register}

\subsection{High Severity}

\begin{enumerate}[leftmargin=*]

\item[\textbf{[C1]}] \textbf{Serial Compute Workgroup in \texttt{vegetation\_instance\_gen.comp}}
  \hfill \severitybox{highcolor}{HIGH}

  \begin{tabular}{ll}
    \textbf{File:} & \texttt{shaders/vegetation\_instance\_gen.comp} \\
    \textbf{Line:} & 3 \\
    \textbf{Category:} & Compute occupancy / GPU utilization \\
  \end{tabular}

  \textbf{Description:}
  The instance-generation compute shader is dispatched with a workgroup of a
  single thread:

\begin{lstlisting}
layout(local_size_x = 1, local_size_y = 1, local_size_z = 1) in;
\end{lstlisting}

  The host dispatch (see \texttt{vulkan/VulkanApp.cpp:457}) issues one workgroup
  \emph{per triangle} — \texttt{vkCmdDispatch(cmd, thisGroups, 1, 1)} where
  \texttt{thisGroups == triCount}. With one invocation per workgroup the GPU
  launches one thread per triangle and never fills a warp/wavefront: on a typical
  GPU each 32/64-lane unit executes a single useful lane and masks the rest.
  Worse, the shader body does no loop over instances, so each thread handles
  exactly one triangle. This is the textbook ``\texttt{local\_size\_x = 1}''
  anti-pattern and is by far the largest single waste in the vegetation path:
  the wavefront is \textasciitilde{}98\% idle.

  \textbf{Recommendation:}
  Raise the workgroup to a hardware-friendly size (e.g. 64) and loop over
  instances inside the shader:

\begin{lstlisting}
layout(local_size_x = 64, local_size_y = 1, local_size_z = 1) in;
void main() {
    uint base = gl_WorkGroupID.x * 64u + gl_LocalInvocationID.x;
    // iterate instances belonging to this triangle, writing into instanceBuffer
}
\end{lstlisting}

  On the CPU side, change the dispatch to
  \texttt{vkCmdDispatch(cmd, (triCount + 63u) / 64u, 1, 1)}. Identical
  instance buffer contents, but the GPU now runs 64× more threads per wavefront.
  No feature or output changes.

\subsection{Medium Severity}

\item[\textbf{[C2]}] \textbf{Per-Invocation Frustum-Plane Normalization in Cull Shaders}
  \hfill \severitybox{mediumcolor}{MEDIUM}

  \begin{tabular}{ll}
    \textbf{File:} & \texttt{shaders/indirect.comp}, \texttt{shaders/vegetation\_cull.comp} \\
    \textbf{Line:} & indirect.comp 37--77; vegetation\_cull.comp 39--69 \\
    \textbf{Category:} & Compute ALU / Invariant Hoisting \\
  \end{tabular}

  \textbf{Description:}
  Both cull shaders reconstruct the four matrix rows from the view-projection
  matrix and then normalize all six frustum planes \emph{inside} the
  \texttt{aabbVisible()} helper, which runs once per invocation:

\begin{lstlisting}
for (int i = 0; i < 6; ++i) {
    vec3 n = planes[i].xyz;
    float d = planes[i].w;
    float len = length(n);      // 6× per invocation
    if (len > 0.0) { n /= len; d /= len; }
    ...
}
\end{lstlisting}

  The view-projection matrix is a \texttt{push\_constant} — identical for every
  invocation of a given dispatch. Six \texttt{length()} calls plus six divides
  are therefore recomputed redundantly for every AABB tested. Across the
  indirect-draw and vegetation-chunk cull passes (potentially tens of thousands
  of invocations per frame) this is pure waste.

  \textbf{Recommendation:}
  Precompute the six normalized planes once on the CPU (or in a tiny one-thread
  setup compute pass) and upload them in a small UBO / push-constant array. The
  shaders then only need the sign-pick + dot test. Alternatively, compute the
  planes once per workgroup using shared memory. Output of the cull test is
  unchanged (the normalization is already what the code does today, just done
  per-invocation).

\item[\textbf{[NL1]}] \textbf{Trigonometric Gradient Hash in \texttt{perlin2d.glsl}}
  \hfill \severitybox{mediumcolor}{MEDIUM}

  \begin{tabular}{ll}
    \textbf{File:} & \texttt{shaders/includes/perlin2d.glsl} \\
    \textbf{Line:} & 46--47 \\
    \textbf{Category:} & Fragment / Vertex ALU \\
  \end{tabular}

  \textbf{Description:}
  Report \#09 (SH3) replaced the \texttt{fract(sin(dot(...)))} hash in
  \texttt{perlin.glsl} and \texttt{perlin4d.glsl} with an integer bit-mixing
  hash. \texttt{perlin2d.glsl} was missed: its gradient direction is still
  derived with transcendental functions:

\begin{lstlisting}
float angle = h * 6.28318530718;
return vec2(cos(angle), sin(angle));   // per corner, per noise sample
\end{lstlisting}

  \texttt{perlin2d} backs the 2D wind/height noise used per-vertex in
  \texttt{vegetation.vert} / \texttt{vegetation.geom} and the impostor density
  test, so \texttt{cos}/\texttt{sin} are issued for every corner of every 2D
  gradient-noise evaluation. This is exactly the same class of problem SH3 fixed
  for the 3D/4D noise, just in the 2D variant.

  \textbf{Recommendation:}
  Mirror the \#09 fix: replace the angle→\texttt{(cos,sin)} gradient with a
  small integer-indexed gradient table (e.g. an 8-direction \texttt{const vec2}
  array selected by \texttt{hash \& 7}), consistent with the
  \texttt{perlin2d\_hash11}/\texttt{hash13} integer hashes already present in
  the same file. Same gradient-noise distribution, no transcendentals in the hot
  path.

\item[\textbf{[T1]}] \textbf{Triplanar Blend Weights Recomputed Three Times per Vertex}
  \hfill \severitybox{mediumcolor}{MEDIUM}

  \begin{tabular}{ll}
    \textbf{File:} & \texttt{shaders/includes/displacement.glsl}, \texttt{shaders/includes/common.glsl} \\
    \textbf{Line:} & displacement.glsl 17/25/33; common.glsl 17--27 \\
    \textbf{Category:} & Tessellation-Eval ALU \\
  \end{tabular}

  \textbf{Description:}
  \texttt{applyDisplacement()} (called once per tessellated vertex in
  \texttt{main.tese}) samples the height of up to three blended materials:

\begin{lstlisting}
h0 = sampleHeightTriplanar(worldPos, worldNormal, texIndices.x);
h1 = sampleHeightTriplanar(worldPos, worldNormal, texIndices.y);
h2 = sampleHeightTriplanar(worldPos, worldNormal, texIndices.z);
\end{lstlisting}

  Each \texttt{sampleHeightTriplanar} call independently rebuilds the triplanar
  blend weights from the \emph{same} \texttt{worldNormal}:

\begin{lstlisting}
vec3 w = abs(normal);                       // common.glsl:19
vec3 wt = max(vec3(0.0), w - vec3(t));     // threshold
wt = pow(wt, vec3(e));                      // 3× pow
float sum = wt.x + wt.y + wt.z + 1e-6;
w = wt / sum;                               // normalize
\end{lstlisting}

  So the identical \texttt{abs} + threshold + \texttt{pow} + normalize sequence
  runs three times per vertex, despite the inputs being constant across the
  three calls. With tessellation enabled this multiplies across every generated
  vertex of every terrain patch.

  \textbf{Recommendation:}
  Compute the triplanar weights \texttt{w} once in \texttt{applyDisplacement()}
  and pass them into a \texttt{sampleHeightTriplanar(worldPos, w, brushIndex)}
  variant that consumes the precomputed weights instead of rebuilding them.
  Identical displaced geometry; one third of the weight math.

\item[\textbf{[GS1]}] \textbf{Fibonacci View-Selection Table Rebuilt per Primitive}
  \hfill \severitybox{mediumcolor}{MEDIUM}

  \begin{tabular}{ll}
    \textbf{File:} & \texttt{shaders/impostors.geom}, \texttt{shaders/impostors\_depth.geom}, \texttt{shaders/vegetation.geom} \\
    \textbf{Line:} & impostors.geom 67--95; impostors\_depth.geom 54--78; vegetation.geom (Fibonacci loop) \\
    \textbf{Category:} & Geometry-shader ALU \\
  \end{tabular}

  \textbf{Description:}
  All three impostor-generating geometry shaders select the best of 20
  Fibonacci-distributed view directions by recomputing the entire direction
  table \emph{inside} the shader body, once per emitted primitive:

\begin{lstlisting}
const float goldenAngle = 3.14159265358979323846 * (3.0 - 2.2360679774997896);
const int   NUM_VIEWS   = 20;
for (int i = 0; i < NUM_VIEWS; i++) {
    float y     = 1.0 - (float(i) + 0.5) / float(NUM_VIEWS) * 2.0;
    float r     = sqrt(max(0.0, 1.0 - y * y));
    float theta = goldenAngle * float(i);
    vec3  d     = vec3(cos(theta) * r, y, sin(theta) * r);  // sqrt + cos + sin ×20
    ...
}
\end{lstlisting}

  The table depends only on the two \texttt{const} values \texttt{goldenAngle}
  and \texttt{NUM\_VIEWS}, so every entry is identical across all invocations
  and across all three shaders. With thousands of impostors per frame this is
  20×(\texttt{sqrt}+\texttt{cos}+\texttt{sin}) repeated per primitive.

  \textbf{Recommendation:}
  Hoist the 20 directions into a compile-time \texttt{const vec3 fibDirs[20]}
  initialized from a precomputed literal, or upload them once as a UBO and index
  by \texttt{i}. The per-primitive loop then becomes a pure dot-product search
  with no transcendentals. Visual selection is unchanged because the values are
  the same.

\subsection{Low Severity}

\item[\textbf{[NL2]}] \textbf{Residual \texttt{sin}-Based Star Hash in Procedural Sky}
  \hfill \severitybox{lowcolor}{LOW}

  \begin{tabular}{ll}
    \textbf{File:} & \texttt{shaders/sky.frag}, \texttt{shaders/sky\_equirect.frag} \\
    \textbf{Line:} & sky.frag 59; sky\_equirect.frag 80 \\
    \textbf{Category:} & Fragment ALU \\
  \end{tabular}

  \textbf{Description:}
  The procedural star field still uses the classic sine hash that SH3 removed
  from the noise library:

\begin{lstlisting}
float starSeed = fract(sin(dot(dir.xz * 1000.0, vec2(12.9898, 78.233))) * 43758.5453);
\end{lstlisting}

  This is a single \texttt{sin} per sky pixel — cheap, and only on the sky dome
  — but it is the same numerical-redundancy class already fixed elsewhere, and
  it sits in a shader that is otherwise hash-clean.

  \textbf{Recommendation:}
  Replace with the same integer bit-mixing hash used in \texttt{perlin.glsl}
  (e.g. a 2D PCG-style mix of the quantized direction). Output distribution
  changes only in the same benign way SH3's change did for noise; no feature
  regression.

\end{enumerate}

\newpage

\section{Summary}

\begin{center}
\begin{tabular}{cll}
\toprule
\textbf{ID} & \textbf{Shader / Include} & \textbf{Severity} \\
\midrule
C1  & \texttt{vegetation\_instance\_gen.comp}        & High   \\
C2  & \texttt{indirect.comp} / \texttt{vegetation\_cull.comp} & Medium \\
NL1 & \texttt{includes/perlin2d.glsl}               & Medium \\
T1  & \texttt{includes/displacement.glsl} / \texttt{common.glsl} & Medium \\
GS1 & \texttt{impostors.geom} / \texttt{impostors\_depth.geom} / \texttt{vegetation.geom} & Medium \\
NL2 & \texttt{sky.frag} / \texttt{sky\_equirect.frag} & Low \\
\bottomrule
\end{tabular}
\end{center}

All six issues are pure-performance: none alters the rendered image. The
highest-impact fix is \textbf{C1} (serial compute workgroup), which is a
one-line shader change plus a matching dispatch-count change on the CPU and
should be done first. C2, T1 and GS1 are hoisting/merge opportunities that
remove redundant invariant math; NL1 and NL2 finish the
\texttt{sin}/\texttt{cos}-hash cleanup begun in Report \#09.

\end{document}
